\documentclass[a4paper]{book}

% 数学宏包
\usepackage{anyfontsize}
\usepackage{amsmath,amssymb,amsthm,mathrsfs}

% Garamond 风格
\usepackage{fontspec}
\setmainfont{EB Garamond}

% 数学字体
\usepackage[scr=rsfso,frak=euler,bb=ams]{mathalfa}
\usepackage[math-style=ISO, bold-style=ISO]{unicode-math}
\usepackage{bm}
\setmathfont{Garamond-Math.otf}[StylisticSet={6,7,9,10}]

% 数学命令
\AtBeginDocument{
  \let\uglyepsilon\epsilon
  \renewcommand\epsilon{\varepsilon}
}
\AtBeginDocument{
  \let\uglymathsf\mathsf
  \renewcommand\mathsf{\symsf}
}
\AtBeginDocument{
  \let\uglyleq\leq
  \renewcommand\leq{\leqslant}
}
\AtBeginDocument{
  \let\uglygeq\geq
  \renewcommand\geq{\geqslant}
}
\renewcommand{\bm}{\symbf}

% 文章排版
\usepackage{indentfirst}
\setlength{\parindent}{0em}
\usepackage{autobreak}
\allowdisplaybreaks
\usepackage{parskip}
\usepackage{geometry}
\geometry{top=3cm,bottom=3cm,left=2cm,right=2cm}

% 中文
\usepackage[fontset=none]{ctex}
\setCJKmainfont{SimSun}[BoldFont=Noto Sans CJK SC Medium]
\setCJKmonofont{Noto Sans Mono CJK SC}[BoldFont=Noto Sans Mono CJK SC Bold]

% 代码
\setmonofont{JetBrains Mono NL}
\usepackage{color}
\usepackage{listings}
\lstset{
    language=Python,
    basicstyle=\small\ttfamily,
    aboveskip={1.0\baselineskip},
    belowskip={1.0\baselineskip},
    columns=fixed,
    extendedchars=true,
    breaklines=true,
    tabsize=4,
    frame=lines,
    showtabs=false,
    showspaces=false,
    showstringspaces=false,
    keywordstyle=\color[rgb]{0,0.5,0},
    commentstyle=\color[rgb]{0.5,0,0},
    stringstyle=\color[rgb]{1,0,0},
    numbers=left,
    numberstyle=\small\ttfamily ,
    stepnumber=1,
    numbersep=12pt,
    captionpos=t,
    escapeinside={\%*}{*)}
}

% 图表
\usepackage{graphicx,caption,subcaption}
\usepackage{supertabular,multirow}

\begin{document}

\title{\textbf{计算物理学作业}}
\author{1900011604\ \ 张植竣\ \ 北京大学}
\date{2021年11月13日}
\maketitle

\begin{enumerate}

\item[1.(a)]

使用Neville算法构造内插多项式：

\lstinputlisting{Neville_1.py}

给出内插多项式为：
\[
    P_4(x) = 0.99631698x + 0.01995143x^2 - 0.20358546x^3 + 0.02871423x^4
\]
内插效果如图：

\begin{figure}[!ht]
    \centering
    \includegraphics[width=0.75\linewidth]{Neville.png}
    \caption{使用Neville算法进行多项式内插}
\end{figure}

通过计算等距的10001个节点$\xi$的误差$\delta=|f(\xi)-P_n(\xi)|$，估计误差最大为：
\[
    \delta_{max} = 0.00021533
\]
而估计任意节点选择下，最大可能的误差为：
\[
    \Delta = \max|f(x)-P_n(x)| \leq \frac{|x_n-x_0|^n}{(n+1)!}\max_{x_0 \leq \xi \leq x_n}\left|f^{(n+1)}(\xi)\right|
\]
即：
\[
    \Delta_{max} = \frac{\left|\frac{\pi}{2}-0\right|^4}{(4+1)!}\max_{0 \leq \xi \leq 1}\left|\frac{\mathrm{d}^{(4+1)}}{\mathrm{d}x^{(4+1)}}\sin(\xi)\right| = 0.05073390
\]
可见有$\delta_{max}<\Delta_{max}$

\item[1.(b)]

选择节点为：
\[
    x_0 = 0,\ x_1 = 0.1,\ x_2 = 0.2,\ x_3 = 0.3,\ x_4 = \frac{\pi}{2}
\]
给出内插多项式为：
\[
    P_4(x) = 0.99992807x + 0.00136841x^2 - 0.17477156x^3 + 0.01697043x^4
\]
通过计算等距的10001个节点$\xi$的误差$\delta=|f(\xi)-P_n(\xi)|$，估计误差最大为：
\[
    \delta_{max} = 0.00316623 < \Delta_{max}
\]
仍小于各种可能的误差最大值。

\item[1.(c)]

由题意，需要找到最小的$n$使其满足
\[
    \Delta_{max}(n) = \frac{1}{(n+1)!}\left(\frac{\pi}{2}\right)^n < 10^{-8}
\]
因为：（$\Delta_{max}(n)$显然是一个单调递减数列）
\[
    \Delta_{max}(12) = 3.62374983\times10^{-8} > 10^{-8} > \Delta_{max}(13) = 4.06583780\times10^{-9}
\]
故满足条件的最小$n$值为13，即至少需要14个节点（包括头尾）进行插值。

\newpage
\item[2.(a)]

同样使用Neville算法，构造内插多项式为：
\begin{align*}
    P_{20}(x) &= 1 + 7.21644966\times10^{-16}x - 2.41434796\times10^{1}x^2 - 6.21724894\times10^{-14}x^3 + 4.70846227\times10^{2}x^4 \\
    &\quad + 8.81072992\times10^{-13}x^5 - 6.11921995\times10^{3}x^6 - 9.77706804\times10^{-12}x^7 + 4.93175429\times10^{4}x^8 \\
    &\quad + 5.09317033\times10^{-11}x^9 - 2.45249284\times10^{5}x^{10} + -2.03726813\times10^{-10}x^{11} + 7.57287093\times10^{5}x^{12} \\
    &\quad + 6.69388101\times10^{-10}x^{13} - 1.44294496\times10^{6}x^{14} - 8.73114914\times10^{-10}x^{15} + 1.63917741\times10^{6}x^{16} \\
    &\quad + 6.98491931\times10^{-10}x^{17} - 1.01209487\times10^{6}x^{18} - 3.05590220\times10^{-10}x^{19} + 2.60178631\times10^{5}x^{20}
\end{align*}
内插效果如图：

\begin{figure}[!ht]
    \centering
    \includegraphics[width=0.75\linewidth]{Neville_Runge.png}
    \caption{使用Neville算法进行多项式内插，探究Runge现象}
\end{figure}

每个节点和区间中点的$x$，$f(x)$，$P_{20}(x)$以及$|f(x)-P_{20}(x)|$列表如下：

\bigskip
\begin{center}
    \tablefirsthead{
    \hline
    $x$ & $f(x)$ & $P_{20}(x)$ & $|f(x)-P_{20}(x)|$ & $x$ & $f(x)$ & $P_{20}(x)$ & $|f(x)-P_{20}(x)|$ \\
    \hline}
    \tablehead{
    \hline
    $x$ & $f(x)$ & $P_{20}(x)$ & $|f(x)-P_{20}(x)|$ & $x$ & $f(x)$ & $P_{20}(x)$ & $|f(x)-P_{20}(x)|$ \\
    \hline}
    \tabletail{\hline}
    \tablelasttail{\hline}
    \begin{supertabular}{|cccc|cccc|}
        -1.00 & 0.0384615 & 0.0384615 & $1.38803\times10^{-10}$ & -0.95 & 0.0424403 & -39.9524  & 39.9949    \\
        -0.90 & 0.0470588 & 0.0470588 & $9.09923\times10^{-11}$ & -0.85 & 0.052459  & 3.45496   & 3.4025     \\
        -0.80 & 0.0588235 & 0.0588235 & $2.27909\times10^{-11}$ & -0.75 & 0.06639   & -0.447052 & 0.513442   \\
        -0.70 & 0.0754717 & 0.0754717 & $4.94807\times10^{-12}$ & -0.65 & 0.0864865 & 0.202423  & 0.115936   \\
        -0.60 & 0.1       & 0.1       & $4.97741\times10^{-13}$ & -0.55 & 0.116788  & 0.08066   & 0.0361283  \\
        -0.50 & 0.137931  & 0.137931  & $4.84335\times10^{-14}$ & -0.45 & 0.164948  & 0.179763  & 0.0148142  \\
        -0.40 & 0.2       & 0.2       & $3.69149\times10^{-15}$ & -0.35 & 0.246154  & 0.238446  & 0.00770791 \\
        -0.30 & 0.307692  & 0.307692  & $2.22045\times10^{-16}$ & -0.25 & 0.390244  & 0.395093  & 0.00484915 \\
        -0.20 & 0.5       & 0.5       & $3.33067\times10^{-16}$ & -0.15 & 0.64      & 0.636755  & 0.00324466 \\
        -0.10 & 0.8       & 0.8       & $1.11022\times10^{-16}$ & -0.05 & 0.941176  & 0.94249   & 0.00131391 \\
        0.00  & 1         & 1         & 0                       & 0.05  & 0.941176  & 0.94249   & 0.00131391 \\
        0.10  & 0.8       & 0.8       & 0                       & 0.15  & 0.64      & 0.636755  & 0.00324466 \\
        0.20  & 0.5       & 0.5       & $1.66533\times10^{-16}$ & 0.25  & 0.390244  & 0.395093  & 0.00484915 \\
        0.30  & 0.307692  & 0.307692  & $1.99840\times10^{-15}$ & 0.35  & 0.246154  & 0.238446  & 0.00770791 \\
        0.40  & 0.2       & 0.2       & $3.63598\times10^{-15}$ & 0.45  & 0.164948  & 0.179763  & 0.0148142  \\
        0.50  & 0.137931  & 0.137931  & $2.47302\times10^{-14}$ & 0.55  & 0.116788  & 0.08066   & 0.0361283  \\
        0.60  & 0.1       & 0.1       & $2.13496\times10^{-13}$ & 0.65  & 0.0864865 & 0.202423  & 0.115936   \\
        0.70  & 0.0754717 & 0.0754717 & $4.52111\times10^{-13}$ & 0.75  & 0.06639   & -0.447052 & 0.513442   \\
        0.80  & 0.0588235 & 0.0588235 & $1.00579\times10^{-11}$ & 0.85  & 0.052459  & 3.45496   & 3.4025     \\
        0.90  & 0.0470588 & 0.0470588 & $1.09567\times10^{-11}$ & 0.95  & 0.0424403 & -39.9524  & 39.9949    \\
        1.00  & 0.0384615 & 0.0384615 & $2.23876\times10^{-11}$ & -     & -         & -         & -          \\
    \end{supertabular}
\end{center}
\bigskip

可见在$0.7<|x|<1$的范围内，误差$|f(x)-P_{20}(x)|$急剧增大，最大误差为：
\[
    \delta_{max} = 39.9948893515
\]
这甚至远远超过了$f(x)$的最大值$f_{max}=1$。

\item[2.(b)]

取阶数$N=20$的标准Chebyshev节点如下：
\[
    x_k = \cos\left(\frac{\pi(k+1/2)}{20}\right),\quad k=0,1,\cdots,19
\]
分别用公式：
\[
    c_m = \frac{2}{N}\sum_{k=0}^{N-1} f\left[\cos\left(\frac{\pi(k+1/2)}{20}\right)\right]\cos\left(\frac{\pi m(k+1/2)}{20}\right)
\]
以及Clenshaw迭代算法，构造目标函数的近似多项式：
\[
    f(x) \approx \frac{c_0}{2} + \sum_{k=1}^{N-1} c_k T_k(x) = \sum_{k=1}^{N} s_k x^{k}
\]
代码如下：

\lstinputlisting{Chebyshev.py}

系数矩阵分别为：
\[
    C =
    \begin{bmatrix}
        3.91955054\times10^{-1} &  9.97465999\times10^{-18} & -2.63311458\times10^{-1} &  1.12757026\times10^{-17} \\
        1.76797696\times10^{-1} & -1.16660154\times10^{-16} & -1.18571565\times10^{-1} &  1.55691432\times10^{-16} \\
        7.93168838\times10^{-2} &  1.56125113\times10^{-17} & -5.27529045\times10^{-2} & -5.37764278\times10^{-17} \\
        3.46293898\times10^{-2} & -4.46691295\times10^{-17} & -2.20465775\times10^{-2} &  2.15756232\times10^{-16} \\
        1.29912177\times10^{-2} & -3.81639165\times10^{-17} & -6.01445262\times10^{-3} & -2.86229374\times10^{-17} \\
    \end{bmatrix}
\]
\[
    S = 
    \begin{bmatrix}
        9.62409671\times10^{-1} & -4.88693287\times10^{-15} & -1.65421760\times10^{1} &  1.32987371\times10^{-13} \\
        1.65458229\times10^{2}  & -1.24297100\times10^{-12} & -9.60824747\times10^{2} &  4.65544270\times10^{-12} \\
        3.37901728\times10^{3}  & -5.00327557\times10^{-12} & -7.41345341\times10^{3} & -1.29798394\times10^{-11} \\
        1.01954749\times10^{4}  &  4.59232652\times10^{-11} & -8.53489389\times10^{3} & -5.71169778\times10^{-11} \\
        3.97316472\times10^{3}  &  3.31397132\times10^{-11} & -7.88326334\times10^{2} & -7.50333129\times10^{-12} \\
    \end{bmatrix}
\]
内插效果如图：

\begin{figure}[ht!]
    \centering
    \begin{subfigure}[t]{0.4\linewidth}
        \centering
        \includegraphics[width=1\linewidth]{Chebyshev.png}
        \subcaption{直接计算Chebyshev多项式的系数}
    \end{subfigure}
    \quad
    \begin{subfigure}[t]{0.4\linewidth}
        \centering
        \includegraphics[width=1\linewidth]{Cheb_Clen.png}
        \subcaption{使用Clenshaw迭代算法}
    \end{subfigure}
    \caption{使用Chebyshev多项式进行内插}
\end{figure}

两种计算方法给出的内插多项式在精度允许范围内一致，通过计算等距的10001个节点$\xi$的误差$\delta=|f(\xi)-C_n(\xi)|$估计最大误差均为：
\[
    \delta_{max} = 0.0375903289
\]
每个节点和区间中点的$x$，$f(x)$，$C_{20}(x)$以及$|f(x)-C_{20}(x)|$列表如下：

\bigskip
\begin{center}
    \tablefirsthead{
    \hline
    $x$ & $f(x)$ & $C_{20}(x)$ & $|f(x)-C_{20}(x)|$ \\
    \hline}
    \tablehead{
    \hline
    $x$ & $f(x)$ & $C_{20}(x)$ & $|f(x)-C_{20}(x)|$ \\
    \hline}
    \tabletail{\hline}
    \tablelasttail{\hline}
    \begin{supertabular}{|cccc|}
        $-9.96917\times10^{-1}$ & $3.86906\times10^{-2}$ & $3.86906\times10^{-2}$ & $1.73953\times10^{-12}$ \\
        $-9.84644\times10^{-1}$ & $3.96227\times10^{-2}$ & $4.10124\times10^{-2}$ & $1.38975\times10^{-3}$  \\
        $-9.72370\times10^{-1}$ & $4.05884\times10^{-2}$ & $4.05884\times10^{-2}$ & $9.80646\times10^{-13}$ \\
        $-9.48125\times10^{-1}$ & $4.26012\times10^{-2}$ & $4.10278\times10^{-2}$ & $1.57346\times10^{-3}$  \\
        $-9.23880\times10^{-1}$ & $4.47651\times10^{-2}$ & $4.47651\times10^{-2}$ & $2.25382\times10^{-3}$  \\
        $-8.88260\times10^{-1}$ & $4.82506\times10^{-2}$ & $5.00512\times10^{-2}$ & $1.80065\times10^{-3}$  \\
        $-8.52640\times10^{-1}$ & $5.21516\times10^{-2}$ & $5.21516\times10^{-2}$ & $1.78940\times10^{-13}$ \\
        $-8.06523\times10^{-1}$ & $5.79308\times10^{-2}$ & $5.57609\times10^{-2}$ & $2.16985\times10^{-3}$  \\
        $-7.60406\times10^{-1}$ & $6.47022\times10^{-2}$ & $6.47022\times10^{-2}$ & $1.11189\times10^{-13}$ \\
        $-7.04927\times10^{-1}$ & $7.44987\times10^{-2}$ & $7.72938\times10^{-2}$ & $2.79513\times10^{-3}$  \\
        $-6.49448\times10^{-1}$ & $8.66208\times10^{-2}$ & $8.66208\times10^{-2}$ & $3.45557\times10^{-15}$ \\
        $-5.85973\times10^{-1}$ & $1.04339\times10^{-1}$ & $1.00421\times10^{-1}$ & $3.91822\times10^{-3}$  \\
        $-5.22499\times10^{-1}$ & $1.27794\times10^{-1}$ & $1.27794\times10^{-1}$ & $2.16493\times10^{-15}$ \\
        $-4.52591\times10^{-1}$ & $1.63373\times10^{-1}$ & $1.69511\times10^{-1}$ & $6.13821\times10^{-3}$  \\
        $-3.82683\times10^{-1}$ & $2.14539\times10^{-1}$ & $2.14539\times10^{-1}$ & $1.94289\times10^{-16}$ \\
        $-3.08064\times10^{-1}$ & $2.96508\times10^{-1}$ & $2.85364\times10^{-1}$ & $1.11436\times10^{-2}$  \\
        $-2.33445\times10^{-1}$ & $4.23295\times10^{-1}$ & $4.23295\times10^{-1}$ & $7.77156\times10^{-16}$ \\
        $-1.55952\times10^{-1}$ & $6.21880\times10^{-1}$ & $6.45255\times10^{-1}$ & $2.33756\times10^{-2}$  \\
        $-7.84591\times10^{-2}$ & $8.66629\times10^{-1}$ & $8.66629\times10^{-1}$ & $1.11022\times10^{-16}$ \\
        $6.24500\times10^{-7}$  & $1.00000\times10^{+0}$ & $9.62410\times10^{-1}$ & $3.75903\times10^{-2}$  \\
        $7.84591\times10^{-2}$  & $8.66629\times10^{-1}$ & $8.66629\times10^{-1}$ & $4.44089\times10^{-16}$ \\
        $1.55952\times10^{-1}$  & $6.21880\times10^{-1}$ & $6.45255\times10^{-1}$ & $2.33756\times10^{-2}$  \\
        $2.33445\times10^{-1}$  & $4.23295\times10^{-1}$ & $4.23295\times10^{-1}$ & $5.55112\times10^{-16}$ \\
        $3.08064\times10^{-1}$  & $2.96508\times10^{-1}$ & $2.85364\times10^{-1}$ & $1.11436\times10^{-2}$  \\
        $3.82683\times10^{-1}$  & $2.14539\times10^{-1}$ & $2.14539\times10^{-1}$ & $2.77556\times10^{-16}$ \\
        $4.52591\times10^{-1}$  & $1.63373\times10^{-1}$ & $1.69511\times10^{-1}$ & $6.13821\times10^{-3}$  \\
        $5.22499\times10^{-1}$  & $1.27794\times10^{-1}$ & $1.27794\times10^{-1}$ & $1.36002\times10^{-15}$ \\
        $5.85973\times10^{-1}$  & $1.04339\times10^{-1}$ & $1.00421\times10^{-1}$ & $3.91822\times10^{-3}$  \\
        $6.49448\times10^{-1}$  & $8.66208\times10^{-2}$ & $8.66208\times10^{-2}$ & $6.89726\times10^{-15}$ \\
        $7.04927\times10^{-1}$  & $7.44987\times10^{-2}$ & $7.72938\times10^{-2}$ & $2.79513\times10^{-3}$  \\
        $7.60406\times10^{-1}$  & $6.47022\times10^{-2}$ & $6.47022\times10^{-2}$ & $7.46070\times10^{-14}$ \\
        $8.06523\times10^{-1}$  & $5.79308\times10^{-2}$ & $5.57609\times10^{-2}$ & $2.16985\times10^{-3}$  \\
        $8.52640\times10^{-1}$  & $5.21516\times10^{-2}$ & $5.21516\times10^{-2}$ & $2.36894\times10^{-14}$ \\
        $8.88260\times10^{-1}$  & $4.82506\times10^{-2}$ & $5.00512\times10^{-2}$ & $1.80065\times10^{-3}$  \\
        $9.23880\times10^{-1}$  & $4.47651\times10^{-2}$ & $4.47651\times10^{-2}$ & $1.52843\times10^{-13}$ \\
        $9.48125\times10^{-1}$  & $4.26012\times10^{-2}$ & $4.10278\times10^{-2}$ & $1.57346\times10^{-3}$  \\
        $9.72370\times10^{-1}$  & $4.05884\times10^{-2}$ & $4.05884\times10^{-2}$ & $1.83474\times10^{-12}$ \\
        $9.84644\times10^{-1}$  & $3.96227\times10^{-2}$ & $4.10124\times10^{-2}$ & $1.38975\times10^{-3}$  \\
        $9.96917\times10^{-1}$  & $3.86906\times10^{-2}$ & $3.86906\times10^{-2}$ & $2.02421\times10^{-13}$ \\
    \end{supertabular}
\end{center}
\bigskip

与2.(a)中的Neville算法内插比较，可见$x$在靠近边界处的误差明显减小，不会出现极端的误差情况。Chebyshev算法的最大误差出现在$x=0$的点，说明Chebyshev多项式的值域限制$[-1,1]$使得其对变化剧烈的函数内插精度不高。

\item[2.(c)]

使用21个等距节点构造三次样条函数，边界条件为Neumann边界条件：
\[
    f'(-1) = S'_\Delta(Y;-1) = -\frac{25}{338},\quad f'(1) = S'_\Delta(Y;1) = \frac{25}{338}
\]
代码如下：

\lstinputlisting{Spline.py}

内插效果如图：

\begin{figure}[!ht]
    \centering
    \includegraphics[width=0.75\linewidth]{Spline.png}
    \caption{使用三次样条函数进行内插}
\end{figure}

通过计算等距的10001个节点$\xi$的误差$\delta=|f(\xi)-S_\Delta(\xi)|$估计最大误差为：
\[
    \delta_{max} = 0.0031689146
\]
每个节点和区间中点的$x$，$f(x)$，$S_\Delta(x)$以及$|f(x)-S_\Delta(x)|$列表如下：

\bigskip
\begin{center}
    \tablefirsthead{
    \hline
    $x$ & $f(x)$ & $S_\Delta(x)$ & $|f(x)-S_\Delta(x)|$ \\
    \hline}
    \tablehead{
    \hline
    $x$ & $f(x)$ & $S_\Delta(x)$ & $|f(x)-S_\Delta(x)|$ \\
    \hline}
    \tabletail{\hline}
    \tablelasttail{\hline}
    \begin{supertabular}{|cccc|}
        -1.00 & 3.84615385e-02 & 3.84615385e-02 & 0.00000000e+00 \\
        -0.95 & 4.24403183e-02 & 4.16088317e-02 & 8.31486638e-04 \\
        -0.90 & 4.70588235e-02 & 4.70588235e-02 & 2.77555756e-16 \\
        -0.85 & 5.24590164e-02 & 5.26792435e-02 & 2.20227074e-04 \\
        -0.80 & 5.88235294e-02 & 5.88235294e-02 & 8.32667268e-17 \\
        -0.75 & 6.63900415e-02 & 6.63276137e-02 & 6.24277479e-05 \\
        -0.70 & 7.54716981e-02 & 7.54716981e-02 & 9.71445147e-17 \\
        -0.65 & 8.64864865e-02 & 8.64914336e-02 & 4.94713838e-06 \\
        -0.60 & 1.00000000e-01 & 1.00000000e-01 & 2.49800181e-16 \\
        -0.55 & 1.16788321e-01 & 1.16782104e-01 & 6.21685096e-06 \\
        -0.50 & 1.37931034e-01 & 1.37931034e-01 & 1.94289029e-16 \\
        -0.45 & 1.64948454e-01 & 1.64865836e-01 & 8.26180987e-05 \\
        -0.40 & 2.00000000e-01 & 2.00000000e-01 & 1.11022302e-16 \\
        -0.35 & 2.46153846e-01 & 2.46267816e-01 & 1.13970091e-04 \\
        -0.30 & 3.07692308e-01 & 3.07692308e-01 & 2.77555756e-16 \\
        -0.25 & 3.90243902e-01 & 3.89419663e-01 & 8.24239000e-04 \\
        -0.20 & 5.00000000e-01 & 5.00000000e-01 & 1.66533454e-16 \\
        -0.15 & 6.40000000e-01 & 6.43168915e-01 & 3.16891462e-03 \\
        -0.10 & 8.00000000e-01 & 8.00000000e-01 & 0.00000000e+00 \\
        -0.05 & 9.41176471e-01 & 9.38866217e-01 & 2.31025403e-03 \\
         0.00 & 1.00000000e+00 & 1.00000000e+00 & 0.00000000e+00 \\
         0.05 & 9.41176471e-01 & 9.38866219e-01 & 2.31025143e-03 \\
         0.10 & 8.00000000e-01 & 8.00000000e-01 & 2.22044605e-16 \\
         0.15 & 6.40000000e-01 & 6.43168907e-01 & 3.16890683e-03 \\
         0.20 & 5.00000000e-01 & 5.00000000e-01 & 1.11022302e-16 \\
         0.25 & 3.90243902e-01 & 3.89419692e-01 & 8.24210434e-04 \\
         0.30 & 3.07692308e-01 & 3.07692308e-01 & 2.22044605e-16 \\
         0.35 & 2.46153846e-01 & 2.46267710e-01 & 1.13863617e-04 \\
         0.40 & 2.00000000e-01 & 2.00000000e-01 & 5.55111512e-16 \\
         0.45 & 1.64948454e-01 & 1.64866233e-01 & 8.22207703e-05 \\
         0.50 & 1.37931034e-01 & 1.37931034e-01 & 1.11022302e-16 \\
         0.55 & 1.16788321e-01 & 1.16780621e-01 & 7.69969104e-06 \\
         0.60 & 1.00000000e-01 & 1.00000000e-01 & 4.44089210e-16 \\
         0.65 & 8.64864865e-02 & 8.64969677e-02 & 1.04811703e-05 \\
         0.70 & 7.54716981e-02 & 7.54716981e-02 & 2.77555756e-17 \\
         0.75 & 6.63900415e-02 & 6.63069605e-02 & 8.30810353e-05 \\
         0.80 & 5.88235294e-02 & 5.88235294e-02 & 1.87350135e-16 \\
         0.85 & 5.24590164e-02 & 5.27563226e-02 & 2.97306191e-04 \\
         0.90 & 4.70588235e-02 & 4.70588235e-02 & 5.96744876e-16 \\
         0.95 & 4.24403183e-02 & 4.13211685e-02 & 1.11914982e-03 \\
         1.00 & 3.84615385e-02 & 3.84615385e-02 & 2.80331314e-15 \\
    \end{supertabular}
\end{center}
\bigskip

与Neville算法和Chebyshev算法比较可见三次样条函数在变化剧烈的$x=0$附近和边界$|x|=1$附近都有较好的内插效果，误差显著减小。

\newpage
\item[3.(a)]

$x(t)$与$y(t)$的精确数值列表如下：

\begin{table}[ht!]
    \begin{center}
    \begin{tabular}{|ccc|}
        \hline
        $t$ & $x$ & $y$ \\
        \hline
        0.00000000e+00 &  0.00000000e+00 & 0.00000000e+00 \\
        3.92699082e-01 &  7.03261419e-02 & 2.91300418e-02 \\
        7.85398163e-01 &  2.07106781e-01 & 2.07106781e-01 \\
        1.17809725e+00 &  2.36236823e-01 & 5.70326142e-01 \\
        1.57079633e+00 &  0.00000000e+00 & 1.00000000e+00 \\
        1.96349541e+00 & -5.29130042e-01 & 1.27743292e+00 \\
        2.35619449e+00 & -1.20710678e+00 & 1.20710678e+00 \\
        2.74889357e+00 & -1.77743292e+00 & 7.36236823e-01 \\
        3.14159265e+00 & -2.00000000e+00 & 2.44929360e-16 \\
        \hline
    \end{tabular}
\end{center}
\end{table}

\item[3.(b)]

过这些点的三次样条函数如下：

\begin{table}[ht!]
    \begin{center}
    \begin{tabular}{|cccc|}
        \hline
        \multicolumn{4}{|c|}{$S_\Delta(X;t_i)=s_0+s_1 t^1+s_2 t^2+s_3 t^3$} \\
        \hline
        $s_0$ & $s_1$ & $s_2$ & $s_3$ \\
        \hline
        5.00000000e-01 &  3.12709949e-01 &  0.00000000 & -7.94575394e-02 \\
        6.83243333e-01 & -3.87228023e-01 &  8.91188705 & -4.57689773e-01 \\
       -3.80796031e+00 &  8.19033909e+00 & -4.56946012 &  7.01095898e-01 \\
       -2.12428116e+00 &  6.04661222e+00 & -3.65963418 &  5.72381857e-01 \\
        3.33705793e+01 & -2.78484827e+01 &  7.12950964 & -5.72381857e-01 \\
        4.11653901e+01 & -3.38032796e+01 &  8.64588621 & -7.01095898e-01 \\
       -8.00971081e+01 &  4.33948244e+01 & -7.73606027 &  4.57689773e-01 \\
       -1.72446451e+01 &  9.09786382e+00 & -1.49773933 &  7.94575394e-02 \\
        \hline
    \end{tabular}
\end{center}
\end{table}

\begin{table}[ht!]
    \begin{center}
    \begin{tabular}{|cccc|}
        \hline
        \multicolumn{4}{|c|}{$S_\Delta(Y;t_i)=s_0+s_1 t^1+s_2 t^2+s_3 t^3$} \\
        \hline
        $s_0$ & $s_1$ & $s_2$ & $s_3$ \\
        \hline
        0.00000000e+00 &  4.27956500e-02 &  0.00000000 &  3.58111040e-01 \\
        4.55184500e-01 & -1.69588107e+00 &  2.21375195 & -5.81434469e-01 \\
       -5.18450015e-02 & -7.27526050e-01 &  1.59727800 & -4.50614633e-01 \\
       -1.47618635e+01 &  1.80018512e+01 & -6.35171659 &  6.73938061e-01 \\
       -1.47618635e+01 &  1.80018512e+01 & -6.35171659 &  6.73938061e-01 \\
        5.33400740e+01 & -3.40241967e+01 &  6.89660773 & -4.50614633e-01 \\
        6.70298705e+01 & -4.27393919e+01 &  8.74602958 & -5.81434469e-01 \\
       -8.90984130e+01 &  4.24557673e+01 & -6.75023408 &  3.58111040e-01 \\
        \hline
    \end{tabular}
\end{center}
\end{table}

\newpage
内插效果如图：

\begin{figure}[ht!]
    \centering
    \begin{subfigure}[t]{0.4\linewidth}
        \centering
        \includegraphics[width=1\linewidth]{Cardioid_x.png}
        \subcaption{$x=(1-\cos{t})\cos{t}$的样条函数内插}
    \end{subfigure}
    \quad
    \begin{subfigure}[t]{0.4\linewidth}
        \centering
        \includegraphics[width=1\linewidth]{Cardioid_y.png}
        \subcaption{$y=(1-\cos{t})\sin{t}$的样条函数内插}
    \end{subfigure}
    \caption{分别对$x-t$，$y-t$进行样条函数内插}
\end{figure}

这里使用的边界条件是：$S''_\Delta(X;0)=S''_\Delta(X;2\pi)=0$，$S''_\Delta(Y;0)=S''_\Delta(Y;2\pi)=0$

\item[3.(c)]

参数形式的曲线如下图：

\begin{figure}[!ht]
    \centering
    \includegraphics[width=0.5\linewidth]{Cardioid_1.png}
    \caption{参数形式的心形线}
\end{figure}

不同边界条件给出的参数形式曲线如下图所示，可见Neumann边界条件和周期性边界条件的效果都不如直接给定端点处二阶导$x''(t)=y''(t)=0$的效果：

\begin{figure}[ht!]
    \centering
    \begin{subfigure}[t]{0.3\linewidth}
        \centering
        \includegraphics[width=1\linewidth]{Cardioid_1.png}
        \subcaption{端点处$x''(t)=y''(t)=0$}
    \end{subfigure}
    \quad
    \begin{subfigure}[t]{0.3\linewidth}
        \centering
        \includegraphics[width=1\linewidth]{Cardioid_2.png}
        \subcaption{端点处$x'(t)=y'(t)=0$}
    \end{subfigure}
    \quad
    \begin{subfigure}[t]{0.3\linewidth}
        \centering
        \includegraphics[width=1\linewidth]{Cardioid_3.png}
        \subcaption{周期性边界条件}
    \end{subfigure}
    \caption{不同边界条件给出的参数形式曲线}
\end{figure}

\item[3.(d)]

因为三次样条函数是使得函数在区间$[a,b]$上的模：
\[
    \|f\| = \int_a^b |f''(x)|^2\,\mathrm{d}x
\]
最小的函数，而这个模可以视为对$f(x)$的曲率$\rho(x)\approx|f''(x)|$的近似。三次样条函数是$f(x)$最小曲率的近似，是“最光滑”的内插函数。它满足在每一个点左右有相同的函数值、一阶导数和二阶导数值，即函数连续、斜率和曲率相同，可以平滑连接所有点。

\newpage
\item[4.(a)]

利用梯形法计算积分，代码如下：

\lstinputlisting{Trapezoidal.py}

\newpage
固定积分步长为$0.1$，选取不同的积分区间，得到的积分值$I$和误差$\delta=\left|\sqrt{\pi}\mathrm{e}^{\tfrac{1}{4}}-I\right|$列表如下：

\begin{table}[ht!]
    \begin{center}
    \begin{tabular}{|ccc|ccc|}
        \hline
        $[-a,a]$ & $I$ & $\delta$ & $[-a,a]$ & $I$ & $\delta$ \\
        \hline
        $[-1,1]$ & 1.31117012 & 6.92183287e-02 & $[-10,10]$ & 1.38038845 & 6.66133815e-16 \\
        $[-2,2]$ & 1.38524609 & 4.85764216e-03 & $[-11,11]$ & 1.38038845 & 4.44089210e-16 \\
        $[-3,3]$ & 1.38042838 & 3.99288253e-05 & $[-12,12]$ & 1.38038845 & 4.44089210e-16 \\
        $[-4,4]$ & 1.38038846 & 1.63180258e-08 & $[-13,13]$ & 1.38038845 & 4.44089210e-16 \\
        $[-5,5]$ & 1.38038845 & 1.05937481e-12 & $[-14,14]$ & 1.38038845 & 4.44089210e-16 \\
        $[-6,6]$ & 1.38038845 & 2.22044605e-16 & $[-15,15]$ & 1.38038845 & 2.22044605e-16 \\
        $[-7,7]$ & 1.38038845 & 2.22044605e-16 & $[-16,16]$ & 1.38038845 & 2.22044605e-16 \\
        $[-8,8]$ & 1.38038845 & 4.44089210e-16 & $[-17,17]$ & 1.38038845 & 2.22044605e-16 \\
        $[-9,9]$ & 1.38038845 & 8.88178420e-16 & $[-18,18]$ & 1.38038845 & 4.44089210e-16 \\
        \hline
    \end{tabular}
\end{center}
\end{table}

可见随着积分区间的增大，误差整体呈现减小的趋势，最后逼近机器精度的极限，误差会有波动。

固定积分区间为$[-10,10]$，选取不同的积分步长$\Delta x=\dfrac{20}{2^N}$，得到的积分值$I$和误差$\delta$列表如下：

\begin{table}[ht!]
    \begin{center}
    \begin{tabular}{|ccc|ccc|}
        \hline
        $N$ & $I$ & $\delta$ & $N$ & $I$ & $\delta$ \\
        \hline
        1 & 10.0000000 & 8.61961155e+00 & 10 & 1.38038845 & 1.55431223e-15 \\
        2 & 5.00000000 & 3.61961155e+00 & 11 & 1.38038845 & 4.44089210e-16 \\
        3 & 2.49226714 & 1.11187870e+00 & 12 & 1.38038845 & 1.77635684e-15 \\
        4 & 1.41136986 & 3.09814175e-02 & 13 & 1.38038845 & 1.33226763e-15 \\
        5 & 1.38038845 & 2.23901231e-09 & 14 & 1.38038845 & 1.24344979e-14 \\
        6 & 1.38038845 & 4.44089210e-16 & 15 & 1.38038845 & 2.39808173e-14 \\
        7 & 1.38038845 & 6.66133815e-16 & 16 & 1.38038845 & 1.35447209e-14 \\
        8 & 1.38038845 & 4.44089210e-16 & 17 & 1.38038845 & 5.70654635e-14 \\
        9 & 1.38038845 & 4.44089210e-16 & 18 & 1.38038845 & 2.23376873e-13 \\
        \hline
    \end{tabular}
\end{center}
\end{table}

可见随着积分区间的增大，误差先减小后增大，区间增大时，函数值相比于机器误差不可忽略，导致总体误差增大。

利用外推积分法计算积分，代码如下：

\lstinputlisting{Extrapolation.py}

选取不同的积分区间，得到的积分值$I$和误差$\delta=\left|\sqrt{\pi}\mathrm{e}^{\tfrac{1}{4}}-I\right|$列表如下：

\begin{table}[ht!]
    \begin{center}
    \begin{tabular}{|ccc|ccc|}
        \hline
        $[-a,a]$ & $I$ & $\delta$ & $[-a,a]$ & $I$ & $\delta$ \\
        \hline
        $[-1,1]$ & 1.31234873 & 6.80397216e-02 & $[-10,10]$ & 1.38038845 & 1.26565425e-14 \\
        $[-2,2]$ & 1.38522290 & 4.83445621e-03 & $[-11,11]$ & 1.38038845 & 3.70814490e-14 \\
        $[-3,3]$ & 1.38042719 & 3.87411007e-05 & $[-12,12]$ & 1.38038845 & 6.12843110e-14 \\
        $[-4,4]$ & 1.38038846 & 1.52084974e-08 & $[-13,13]$ & 1.38038845 & 6.94999613e-14 \\
        $[-5,5]$ & 1.38038845 & 1.01474384e-12 & $[-14,14]$ & 1.38038845 & 8.65973959e-14 \\
        $[-6,6]$ & 1.38038845 & 6.66133815e-16 & $[-15,15]$ & 1.38038845 & 8.05355782e-13 \\
        $[-7,7]$ & 1.38038845 & 6.66133815e-16 & $[-16,16]$ & 1.38038845 & 2.96984659e-12 \\
        $[-8,8]$ & 1.38038845 & 2.66453526e-15 & $[-17,17]$ & 1.38038845 & 7.99982303e-12 \\
        $[-9,9]$ & 1.38038845 & 2.88657986e-15 & $[-18,18]$ & 1.38038845 & 1.78181914e-11 \\
        \hline
    \end{tabular}
\end{center}
\end{table}

可见随着积分步长的减小，误差先减小后增大，步长过小时，机器误差不可忽略且误差累计次数增多，导致总体误差增大。

选取不同的外推阶数$m$，选取积分步长为：
\[
    \Delta x=\dfrac{20}{2^i}, i = 0,1,\cdots,m
\]
得到的积分值$I$和误差$\delta$列表如下：

\begin{table}[ht!]
    \begin{center}
    \begin{tabular}{|ccc|ccc|}
        \hline
        $m$ & $I$ & $\delta$ & $m$ & $I$ & $\delta$ \\
        \hline
        1 & 13.3333333 & 1.19529449e+01 & 10 & 1.38038845 & 1.26565425e-14 \\
        2 & 2.66666667 & 1.28627822e+00 & 11 & 1.38038845 & 2.22044605e-16 \\
        3 & 1.52674646 & 1.46358017e-01 & 12 & 1.38038845 & 2.44249065e-15 \\
        4 & 1.00018790 & 3.80200547e-01 & 13 & 1.38038845 & 1.11022302e-15 \\
        5 & 1.39930803 & 1.89195831e-02 & 14 & 1.38038845 & 1.75415238e-14 \\
        6 & 1.38082752 & 4.39067991e-04 & 15 & 1.38038845 & 2.86437540e-14 \\
        7 & 1.38037481 & 1.36356818e-05 & 16 & 1.38038845 & 8.65973959e-15 \\
        8 & 1.38038851 & 5.99889054e-08 & 17 & 1.38038845 & 7.68274333e-14 \\
        9 & 1.38038845 & 6.02289330e-11 & 18 & 1.38038845 & 2.97317726e-13 \\
        \hline
    \end{tabular}
\end{center}
\end{table}

可见随着外推阶数的增大，误差先减小后增大，阶数过大、步长过小时，机器误差不可忽略且误差累计次数增多，导致总体误差增大。

\newpage
\item[4.(b)]

该积分应该采用Gauss-Hermite积分公式，积分代码如下：

\lstinputlisting{Gauss_Int.py}

其中积分节点和权重系数由函数\texttt{numpy.polynomial.chebyshev.chebgauss(deg)}给出。

选取阶数$deg$分别为5、10、15、20，积分结果$I$和误差$\delta=\left|\sqrt{\pi}\mathrm{e}^{\tfrac{1}{4}}-I\right|$分别为：

\begin{table}[ht!]
    \begin{center}
    \begin{tabular}{|ccc|}
        \hline
        $deg$ & $I$ & $\delta$ \\
        \hline
        5  & 1.3803900759356564 & -1.6288925135388155e-06 \\
        10 & 1.3803884470431410 & 1.9984014443252818e-15  \\
        15 & 1.3803884470431430 & 0.0  \\
        20 & 1.3803884470431427 & 2.220446049250313e-16  \\
        \hline
    \end{tabular}
\end{center}
\end{table}

可见在阶数$deg>10$之后积分结果趋于稳定，误差仅仅受机器精度限制。

\end{enumerate}

\end{document}
